% s-005-p005-m002; original PDF 5, 6, 7
\PrefaceSection{s-005-p005-m002}{2}
Dies war gewiss eine Aufrichtigkeit, und aufrichtig sein dürfen, verführt zum Muthwillen. Die ernsthaften, zugeknöpften, feierlichen Menschen sind immer Poseure, die Einiges sehen lassen und Vieles verbergen; Sterne nennt die Gravität eine Erzschelmin. Indessen, hüten wir uns, dass dieser aufrichtige Muthwille der Vorrede nicht für das Buch ein Versprechen, ein Zuviel-Versprechen bedeute. Man kann zwar mit dem blossen ernsthaften Tonfall Stunden lang predigen und Bände von achthundert Seiten schreiben — aber man kann nicht aus lauter Ausgelassenheit und guter Laune und südlichem Himmel ein Buch machen. Ein wenig Sachlichkeit ist schon nöthig, und so erlaubt sich auch dieses Buch, an Orten, wo man es nicht einmal vermuthen würde, irgend eine vernünftige, bisweilen gar wissenschaftliche Frage zu stellen und auf eigene Faust zu beantworten. Auf diese Weise kommt hier eine Meinung, da eine Einsicht, dort wenigstens die Spitze eines Problems zum Vorschein — und dem Leser bleibt das Vergnügen ungeraubt, diese Einzelheiten zum System in Reih und Glied zu stellen, oder sie als das Chaos zu verdammen. Sollte aber wirklich nicht der Schatten eines Systems darin sein, nicht einmal das bedauerliche System der Systemlosigkeit, sollte es gar — schrecklich zu denken! — nicht am Selbstwiderspruch fehlen, den ich freilich als Recht und als Merkmal des \Sperr{Lebendigen} verehre, nun, so bleibt nichts übrig, als auf eine andere als die systematische Einheit meines Buches zaghaft hinzudeuten. Moderne Menschen lieben das die \Sperr{Einheit der Stimmung} zu nennen, worin freilich liegt, dass man sie mittheilen und empfinden, nicht von ihr sprechen soll. Mein Buch hat im Ganzen einen und denselben Barometerstand; seine Sprüche und Gedanken sind, wenn auch durch Ort und Zeit getrennt, alle in einer mittleren Bergeshöhe ans Licht getreten. Sie entstammen heilig-schattenkühlen Alpensommern und einem sonnenheitren ligurischen Winter, sie erzählen, mehr oder minder unmittelbar, von den wenigen Stätten in Europa, wo es leidliches und bisweilen sogar erfreuliches Wetter giebt. Man wird aus ihnen eine tiefe, fast leidenschaftliche Überzeugung von der Unbewohnbarkeit unserer nordischen Breiten heraushören; die Klimafrage, wichtiger noch als alle Diätfragen, wird ja lange nicht ernst genug behandelt. Was ist die chronische Unvollkommenheit des Nordens und unsere eigentliche Sünde wider den heiligen Geist, wenn nicht das schlechte Wetter! die ewige „kimmerische Halbnacht“, der bleierne unentschiedene Dämmerzustand, der kein Ja oder Nein von Licht und Wärme kennt. Was wissen wir armen Blinden von Farben! Man sollte nie zulassen, dass nordisch blöde Augen mit ihren schleimigen dunstigen Begriffen von Grün und Blau an der Wahrscheinlichkeit Böcklinscher Gemälde herumtasten. Ein deutscher Nebelspalter kann an diese Cypressen, dieses Violett der Schatten, dieses Purpurbraun beschäumter Klippen nicht glauben; diese silbernen und tiefgoldenen Stimmen des Lichts grüssen nur den südwärts Geflüchteten, der morgens aus Zarathustras Höhle in Zarathustras Garten tritt! Im deutschen Klima, das keine Nüancen und Einzelheiten zu erkennen erlaubt, wird man auch „systemsichtig“, man sieht die Welt farblos, vereinfacht, abstract; bei Nacht sind alle Katzen grau und alle Systeme wahr. Mein systemlos übermüthiges Buch überlässt es den Lesern, seinen Mangel an System darauf zurückzuführen, dass es nicht in Deutschland, dem Lande schlechten Wetters und systematischen Denkens, geschrieben wurde. Es überlässt ihnen auch, seinen Mangel an Idealismus zu beklagen: der Deutsche, den es nach dem Süden zieht, legt dieser Sehnsucht einen Bildungsgrund unter — er sucht nicht bloss gutes Wetter, sondern die grosse Kunst der Antike und Renaissance. Aber dieses „nicht bloss“ ist unverzeihlicher Hochmuth; Italiens Kunst ist transformirtes und verewigtes gutes Wetter, wir geniessen in ihr alle heiteren und entwölkten Tage der Vergangenheit!
\par

